% \iffalse meta-comment
%
% Copyright (C) 1993-2026
%
% The LaTeX Project and any individual authors listed elsewhere
% in this file.
%
% This file is part of the Standard LaTeX `Tools Bundle'.
% -------------------------------------------------------
%
% It may be distributed and/or modified under the
% conditions of the LaTeX Project Public License, either version 1.3c
% of this license or (at your option) any later version.
% The latest version of this license is in
%    https://www.latex-project.org/lppl.txt
% and version 1.3c or later is part of all distributions of LaTeX
% version 2008 or later.
%
% The list of all files belonging to the LaTeX `Tools Bundle' is
% given in the file `manifest.txt'.
%
% \fi
%
%\iffalse   % this is a METACOMMENT !
%
%
%% Package `verbatim' to use with LaTeX2e
%% Copyright (C) 1989--2003 by Rainer Sch\"opf. All rights reserved.
%
% Copying of this file is authorized only if either
% (1) you make absolutely no changes to your copy, including name, or
% (2) if you do make changes, you name it something other than
%     verbatim.dtx.
% This restriction helps ensure that all styles developed here
% remain identical.
%
%
%
% \section{Producing the documentation}
%
% We provide a short driver file that can be extracted by the
% \textsf{DocStrip} program using the conditional `\textsf{driver}'.
%
%    \begin{macrocode}
%<*driver>
\documentclass{ltxdoc}
\usepackage[UTF8,scheme=chinese]{ctex}
\pagestyle{myheadings}

\title{\LaTeX{} \texttt{verbatim} 与 \texttt{verbatim*}\\环境的新实现}
\author{Rainer Sch\"opf\\
        \and
        Bernd Raichle\\
        \and
        Chris Rowley}

\date{2023/11/06}
\begin{document}
\markboth{Verbatim 宏包}{Verbatim 宏包}
\MaintainedByLaTeXTeam{tools}
\maketitle
  \DocInput{verbatim-zh.dtx}
\end{document}
%</driver>
%    \end{macrocode}
%
%
%\fi
%
%
% 
% \changes{v1.5x}{2024/01/22}{Added TAB marking support to \cs{verbatiminput*}}
% \changes{v1.5u}{2020-07-07}{Typo repair and added a missing comment
%    character}
% \changes{v1.5q}{2003/08/22}{Reintroduced \cs{@noligs}, by popular
%                             request.}
% \changes{v1.5i}{1996/06/04}{Move setting of verbatim font and
%                             \cs{@noligs}.}
% \changes{v1.5g}{1995/04/26}{Removed \cs{fileversion} and
%                             \cs{filedate} from running head in
%                             driver file, as these are no longer
%                             defined.}
% \changes{v1.5f}{1994/10/25}{Removed extra \cs{typeout} commands.}
% \changes{v1.5e}{1994/06/10}{Added missing closing verbtest guard.}
% \changes{v1.5d}{1994/05/30}{\cs{NeedsTeXFormat} and
%                             \cs{ProvidesPackage} added.}
% \changes{v1.5d}{1994/05/30}{\cs{addto@hook} removed, now in kernel.}
% \changes{v1.5a}{1993/10/12}{Included \cs{newverbtext} command, as
%          written by Chris Rowley.}
% \changes{v1.5}{1993/10/11}{Included vrbinput style option by Bernd
%          Raichle.}
%
% \changes{v1.4j}{1992/06/30}{Used \cs{lowercase}\{\cs{endgroup}
%    \ldots\} trick proposed by Bernd Raichle; changed all \cs{gdef}
%    to \cs{def} since no longer necessary.}
% \changes{v1.4g}{1991/11/21}{Several improvements in the
%                           documentation.}
% \changes{v1.4f}{1991/08/05}{Corrected bug in documentation.
%                           Found by Bernd Raichle.}
% \changes{v1.4e}{1991/07/24}{Avoid reading this file twice.}
% \changes{v1.4d}{1991/04/24}{\cs{penalty}\cs{interlinepenalty} added to
%                           definition of \cs{par} in \cs{@verbatim}.
%                           Necessary to avoid page breaks in
%                           the scope of a \cs{samepage} declaration.}
% \changes{v1.4c}{1990/10/18}{Added \cs{leavevmode} to definition of
%       backquote macro.}
% \changes{v1.4b}{1990/07/14}{Converted nearly all \cs{verb}'s to
%       \texttt{\protect\string!|\ldots\protect\string!|}.}
% \changes{v1.4a}{1990/04/04}{Added a number of percent characters
%       to suppress blank space at the end of some code lines.}
% \changes{v1.4}{1990/03/07}{\cs{verb} rewritten.}
%
% \changes{v1.3a}{1990/02/04}{Removed \texttt{verbatimwrite} environment
%       from the code. Now only shown as an example.}
%
% \changes{v1.2g}{1990/02/01}{Revised documentation.}
% \changes{v1.2e}{1990/01/15}{Added \cs{every@verbatim} hook.}
% \changes{v1.2d}{1989/11/29}{Use token register \cs{@temptokena}
%                           instead of macro \cs{@tempb}.}
% \changes{v1.2d}{1989/11/29}{Use token register \cs{verbatim@line}
%                           instead of macro \cs{@tempd}.}
% \changes{v1.2b}{1989/10/25}{\cs{verbatimfile} renamed to
%           \cs{verbatiminput}. Suggested by Reinhard Wonneberger.}
%
% \changes{v1.1a}{1989/10/16}{\cs{verb} added.}
% \changes{v1.1}{1989/10/09}{Made the code more modular (as suggested by
%                          Chris Rowley):  introduced
%                          \cs{verbatim@addtoline}, etc.  Added
%                          \cs{verbatimwrite} environment.}
%
% \changes{v1.0e}{1989/07/17}{Fixed bug in \cs{verbatimfile} (*-form
%         handling, discovered by Dirk Kreimer).}
% \changes{v1.0d}{1989/05/16}{Revised documentation, fixed silly bug
%         in \cs{verbatim@@@}.}
% \changes{v1.0c}{1989/05/12}{Added redefinition of \cs{@sverb}, change
%         in end-of-line handling.}
% \changes{v1.0b}{1989/05/09}{Change in \cs{verbatim@rescan}.}
% \changes{v1.0a}{1989/05/07}{Change in \cs{verbatim@@testend}.}
%
%
% \DoNotIndex{\ ,\!,\C,\[,\\,\],\^,\`,\{,\},\~}
% \DoNotIndex{\@M,\@empty,\@flushglue,\@gobble,\@ifstar,\@ifundefined}
% \DoNotIndex{\@namedef,\@spaces,\@tempa,\@tempb,\@tempc,\@tempd}
% \DoNotIndex{\@temptokena,\@totalleftmargin,\@warning,\active}
% \DoNotIndex{\aftergroup,\arabic,\begingroup,\catcode,\char,\closein}
% \DoNotIndex{\csname,\def,\do,\docdate,\dospecials,\edef,\else}
% \DoNotIndex{\endcsname,\endgraf,\endgroup,\endinput,\endlinechar}
% \DoNotIndex{\endtrivlist,\expandafter,\fi,\filedate,\fileversion}
% \DoNotIndex{\frenchspacing,\futurelet,\if,\ifcat,\ifeof,\ifnum}
% \DoNotIndex{\ifx,\immediate,\item,\kern,\lccode,\leftskip,\let}
% \DoNotIndex{\lowercase,\m@ne,\makeatletter,\makeatother,\newread}
% \DoNotIndex{\newread,\next,\noexpand,\noindent,\openin,\parfillskip}
% \DoNotIndex{\parindent,\parskip,\penalty,\read,\relax,\rightskip}
% \DoNotIndex{\sloppy,\space,\string,\the,\toks@,\trivlist,\tt,\typeout}
% \DoNotIndex{\vskip,\write,\z@}
%
% \changes{v1.5x}{2025-06-04}{Add notice pointing to \pkg{fancyvrb}}
% \begin{flushleft}
%   \bfseries
%   本宏包出于稳定性考虑仍保留在 \LaTeX{} \pkg{tools} 宏包集合中。
%   虽然会继续修复 \pkg{verbatim} 的 bug，但不再考虑添加\emph{新}功能。
%   对于新的需求，\LaTeX{} 团队建议使用 \pkg{fancyvrb} 代替
%   \pkg{verbatim}。
% \end{flushleft}
%
% \begin{abstract}
%   本宏包重新实现了 \LaTeX{} 的 \texttt{verbatim} 和
%   \texttt{verbatim*} 环境。此外它还提供了一个 \texttt{comment}
%   环境，用于跳过 |\begin{comment}| 与下一个 |\end{comment}|
%   之间的所有命令或文本。它还定义了 |\verbatiminput| 命令，
%   用于以 verbatim 方式包含整个文件。
% \end{abstract}
%
% \section{使用说明}
%
% \let\docDescribeMacro\DescribeMacro
% \let\docDescribeEnv\DescribeEnv
% \def\DescribeMacro#1{}
% \def\DescribeEnv#1{}
% \LaTeX{} 的 \texttt{verbatim} 和 \texttt{verbatim*} 环境
% 有几个可能引发问题的特性：
% \begin{itemize}
%   \item
%     由于检测 |\end{verbatim}| 的方式（即宏参数定界），
%     你不能在 |\end| 记号与 |{verbatim}| 之间留空格。
%   \item
%     由于 \TeX{} 必须先读取 |\begin{verbatim}| 与
%     |\end{verbatim}| 之间的全部文本才能输出任何内容，
%     过长的 verbatim 列表可能会耗尽 \TeX{} 的内存。
% \end{itemize}
% 这两点中，前者只能算小麻烦，后者却是真正的限制。
%
%
% \DescribeEnv{verbatim}
% \DescribeEnv{verbatim*}
% 本宏包文件包含了对 \texttt{verbatim} 和 \texttt{verbatim*}
% 环境的重新实现，克服了上述限制。不过新旧实现之间有一处不兼容：
% 旧版本会把 |\end{verbatim}| 命令同一行上的文本当作单独一行处理。
% \begin{center}
%   \bf 新版本则会直接忽略这些文本。
% \end{center}
% （这是消除旧 \texttt{verbatim} 环境大小限制所必须付出的代价。）
% 不过它会发出如下形式的警告信息：
% \begin{verbatim}
%LaTeX warning: Characters dropped after \end{verbatim*}!
%\end{verbatim}
% 这并不是真正的问题，因为这些文本很容易移到下一行而不影响输出。
%
% 新实现还解决了上面提到的第二个问题：现在可以在 |\end| 与
% |{verbatim}| 或 |{verbatim*}| 之间留空格（但\emph{不能}开始新行）：
% \begin{verbatim}
%\begin {verbatim*}
%   test
%   test
%\end {verbatim*}
%\end{verbatim}
%
% \DescribeEnv{comment}
% 此外我们引入了一个 \texttt{comment} 环境，其效果是：
% |\begin{comment}| 与 |\end{comment}| 之间的文本无论长什么样都被直接忽略。
% 乍看之下这与 verbatim 列表的用途大不相同，但实际上这两个概念的
% 实现非常相似。二者都依赖于这样一个事实：|\begin{...}| 与
% |\end{...}| 之间的文本由 \TeX{} 读取，不解释任何命令或特殊字符。
% \texttt{verbatim} 与 \texttt{comment} 之间剩下的唯一区别是：
% 前者的文本要排版出来，后者的文本则被丢弃。注意这些环境不能嵌套。
%
% \DescribeMacro{\verbatiminput}
% |\verbatiminput| 是一个带一个参数的命令，用于以 verbatim 方式
% 包含一个文件；也就是说，命令 |\verbatiminput{xx.yy}| 的效果等价于\\[2pt]
%   \hspace*{\MacroIndent}|\begin{verbatim}|\\
%   \hspace*{\MacroIndent}\meta{文件 \texttt{xx.yy} 的内容}\\
%   \hspace*{\MacroIndent}|\end{verbatim}|\\[2pt]
% 该命令还有一个 |*|-变体，会把空格打印为 \verb*+ +。
%
%
% \MaybeStop{}
%
%
% \section{面向宏包作者的接口}
%
% \LaTeXe{} 的 \texttt{verbatim} 环境没有为程序员提供良好的接口。
% 与之不同，本宏包提供了一个简单的机制来实现类似功能，
% 这里实现的 \texttt{comment} 环境就是一个示例，展示了能做什么以及怎么做。
%
%
% \subsection{简单示例}
%
% 现在可以使用 \texttt{verbatim} 环境来定义你自己的环境了。
% 例如，
%\begin{verbatim}
% \newenvironment{myverbatim}%
%   {\endgraf\noindent MYVERBATIM:%
%    \endgraf\verbatim}%
%   {\endverbatim}
%\end{verbatim}
% 之后就可以像 \texttt{verbatim} 环境一样使用它，即
% \begin{verbatim}
%\begin {myverbatim}
%  test
%  test
%\end {myverbatim}
%\end{verbatim}
% 另一种用法是写
% \begin{verbatim}
%\let\foo=\comment
%\let\endfoo=\endcomment
%\end{verbatim}
% 此后环境 \texttt{foo} 就与 comment 环境相同，即其内部的所有内容都被忽略。
%
% 你还可以在 |\verbatim| 宏被调用之后添加特殊命令，例如
%\begin{verbatim}
%\newenvironment{myverbatim}%
%  {\verbatim\myspecialverbatimsetup}%
%  {\endverbatim}
%\end{verbatim}
% 不过此时你可能想了解下面要讲的 |\every@verbatim| 钩子。
% \changes{v1.5h}{1995/09/21}{Clarified documentation on use of other
%               environments to define new verbatim-type ones.}
% 不过仍有若干限制：
% \begin{enumerate}
%   \item
%     你不能在定义内部使用 |\begin| 或 |\end| 命令，
%     例如下面两个示例\emph{不能}工作：
%\begin{verbatim*}
%\newenvironment{myverbatim}%
%{\endgraf\noindent MYVERBATIM:%
% \endgraf\begin{verbatim}}%
%{\end{verbatim}}
%\newenvironment{fred}
%{\begin{minipage}{30mm}\verbatim}
%{\endverbatim\end{minipage}}
%\end{verbatim*}
%     如果你尝试这些示例，\TeX{} 会报 ``runaway argument'' 错误。
%     更一般地说，在新环境的定义中不能使用
%     |\begin|\ldots\allowbreak|\end| 或相关环境。正确的定义方式应该是
%    \begin{verbatim*}
%\newenvironment{fred}
%{\minipage{30mm}\verbatim}
%{\endverbatim\endminipage}
%\end{verbatim*}
%   \item
%     你\emph{不能}在用户自定义的\emph{命令}内部使用
%     \texttt{verbatim} 环境；例如
% \changes{v1.4g}{1991/11/21}{Corrected wrong position of optional
%        argument to \cs{newcommand}. Discovered by Piet van Oostrum.}
%     \begin{verbatim*}
%\newcommand{\verbatimfile}[1]%
%           {\begin{verbatim}\input{#1}\end{verbatim}}
%\end{verbatim*}
%     不\emph{能}工作；下面这个也不行：
%     \begin{verbatim}
%\newcommand{\verbatimfile}[1]{\verbatim\input{#1}\endverbatim}
%\end{verbatim}
%   \item 新定义环境的名字不能包含字符类别不是 $11$（字母）
%     或 $12$（其他）的字符，否则将无法工作。
% \end{enumerate}
%
%
% \subsection{接口}
%
% \DescribeMacro{\verbatim@font}
% 从简单的开始讲起。有时你可能需要为 verbatim 文本使用特殊字体，
% 或者使用较小字号的计算机现代打字机字体。
%
% 你可以通过重定义宏 |\verbatim@font| 来选择字体。
% 该宏在每段 verbatim 文本开始时被调用，用于选择字体形状。
% 不要把它用于其他用途；如果你发现自己在滥用它，
% 可以去看看下面的 |\every@verbatim| 钩子。
%
% 默认情况下，|\verbatim@font| 切换到打字机字体并禁用其中的连字。
%
%
% \DescribeMacro{\every@verbatim}
% \DescribeMacro{\addto@hook}
% 有一个钩子（即记号寄存器）叫做 |\every@verbatim|，
% 其内容会在每段 verbatim 文本之前被插入到 \TeX{} 的"嘴"中。
% 请使用 |\addto@hook| 宏向这个钩子添加内容。用法如下：\\[2pt]
% \hspace*{\MacroIndent}|\addto@hook|\meta{钩子名}^^A
%  |{|\meta{要添加的命令}|}|
% \vspace*{2pt}
%
%
%
% \DescribeMacro{\verbatim@start}
% 在完成所有特定设置（如切换字符类别）之后，会调用
% |\verbatim@start| 宏。它启动本宏包实现的扫描机制主循环。
% 任何想利用此功能的其他环境都应把该宏作为最后一个动作来执行。
%
%
% \DescribeMacro{\verbatim@startline}
% \DescribeMacro{\verbatim@addtoline}
% \DescribeMacro{\verbatim@processline}
% \DescribeMacro{\verbatim@finish}
% 这些是与 verbatim 环境启动相关的宏。
% 完成这些（及其他）设置后，本宏包中的代码按以下方式从输入流中
% 读取并处理字符：
% \begin{enumerate}
%   \item 在读取一行输入的第一个字符之前，
%     执行宏 |\verbatim@startline|。
%   \item 读取若干字符后，以这些字符为唯一参数调用宏
%     |\verbatim@addtoline|。每行可能发生多次
%     （当该行含有 |\end| 命令时）。
%   \item 到达行尾时，调用宏 |\verbatim@processline|
%     来处理 |\verbatim@addtoline| 累积的字符。
%   \item 最后，在环境被 |\end| 宏调用结束之前，
%     会调用宏 |\verbatim@finish|。
% \end{enumerate}
%
%
% 为说明这一点，我们看一下标准 \texttt{verbatim} 环境。
% 在这种情况下，上述三个宏定义如下：
% \begin{enumerate}
%   \item |\verbatim@startline| 清空字符缓冲区（一个记号寄存器）。
%   \item |\verbatim@addtoline| 把其参数中的记号
%     添加到字符缓冲区 |\verbatim@line|，不展开它们。
%   \item |\verbatim@processline| 排版缓冲区中累积的字符。
% \end{enumerate}
% 有了这些，实现 \texttt{comment} 环境就非常简单了：
% 在这种情况下，|\verbatim@startline| 和 |\verbatim@processline|
% 被定义为空操作，而 |\verbatim@addtoline| 丢弃其参数。
%
%
% 让我们用这个机制来定义一个 |verbatim| 环境的变体，
% 在左边距打印行号。假设用一个叫 |VerbatimLineNo| 的计数器来实现。
% 假设该计数器已由环境正确初始化，那么在这种情况下
% |\verbatim@processline| 定义为：
% \begin{verbatim}
%\def\verbatim@processline{%
%  \addtocounter{VerbatimLineNo}{1}%
%  \leavevmode
%  \llap{\theVerbatimLineNo\ \hskip\@totalleftmargin}%
%  \the\verbatim@line\par}
%\end{verbatim}
%
% 另一种可能是定义一个 |verbatim| 环境的变体，
% 把整段 verbatim 文本放入盒子并居中。
% 注意盒子中的文本应少于一页，否则你需要修改这个示例。
%
%\begin{verbatim}
%\def\verbatimboxed#1{\begingroup
%  \def\verbatim@processline{%
%    {\setbox0=\hbox{\the\verbatim@line}%
%     \hsize=\wd0
%     \the\verbatim@line\par}}%
%  \setbox0=\vbox{\parskip=0pt\topsep=0pt\partopsep=0pt
%                 \verbatiminput{#1}}%
%  \begin{center}\fbox{\box0}\end{center}%
% \endgroup}
%\end{verbatim}
%
% 最后一个稍有分量的示例是描述一个叫 \texttt{verbatimwrite} 的环境的定义。
% 它把其内部所有文本写入一个文件，文件名作为参数给出。
% 我们假设有一个叫 |\verbatim@out| 的流号已经通过 |\newwrite| 宏预留。
%
% 从宏 |\verbatimwrite| 的定义开始。
% \begin{verbatim}
%\def\verbatimwrite#1{%
%\end{verbatim}
% 首先调用 |\@bsphack|，使该环境不影响间距。
% 然后打开文件并设置所有特殊字符的类别码：
% \changes{v1.5t}{2020-07-05}{Added quotes around the filename in
%    order to allow a filename with spaces; also added a space and
%    comment-character to allow for thespace dlimited argument.}
% \begin{verbatim}
%  \@bsphack
%  \immediate\openout \verbatim@out "#1" %
%  \let\do\@makeother\dospecials
%  \catcode`\^^M\active
%\end{verbatim}
% 下面几个宏的默认定义
% \begin{verbatim}
%  \verbatim@startline
%  \verbatim@addtoline
%  \verbatim@finish
%\end{verbatim}
% 也在这个环境中使用。只需要在调用 |\verbatim@start| 之前
% 修改宏 |\verbatim@processline|：
% \begin{verbatim}
%  \def\verbatim@processline{%
%    \immediate\write\verbatim@out{\the\verbatim@line}}%
%  \verbatim@start}
%\end{verbatim}
% |\endverbatimwrite| 的定义非常简单：
% 关闭流并调用 |\@esphack| 来恢复正确间距。
% \begin{verbatim}
%\def\endverbatimwrite{\immediate\closeout\verbatim@out\@esphack}
%\end{verbatim}
%
% \section{实现}
%
% \let\DescribeMacro\docDescribeMacro
% \let\DescribeEnv\docDescribeEnv
%
% \changes{v1.4e}{1991/07/24}{Avoid reading this file twice.}
% 我们做的第一件事是确保本文件不会被读取两次。为此我们检查
% 宏 |\verbatim@@@| 是否已定义。如果已定义，就停止读取本文件。
% 这里的"宏包"守卫使得在抽取用于测试本宏包的 driver 文件时，
% 大部分代码可以被排除。
%    \begin{macrocode}
%<*package>
\NeedsTeXFormat{LaTeX2e}
\ProvidesPackage{verbatim}
     [2026-08-17 v1.5z LaTeX2e package for verbatim enhancements]
\@ifundefined{verbatim@@@}{}{\endinput}
%    \end{macrocode}
%
% 我们使用与 \AmSTeX 中 |\comment|\ldots\allowbreak|\endcomment|
% 宏类似的机制：一次读取一行，检查它是否包含 |\end{...}| 记号。
% 然后就可以判断是否到达了 verbatim 文本的末尾，还是需要继续。
%
%
% \subsection{预备工作}
%
% \begin{macro}{\every@verbatim}
%    钩子（即记号寄存器）|\every@verbatim| 初始化为 \meta{空}。
%    \begin{macrocode}
\newtoks\every@verbatim
\every@verbatim={}
%    \end{macrocode}
% \end{macro}
%
%
% \begin{macro}{\@makeother}
% \changes{v1.1a}{1989/10/16}{\cs{@makeother} added.}
%    |\@makeother| 以一个字符为参数，把它的类别码改为 $12$（其他）。
%    \begin{macrocode}
\def\@makeother#1{\catcode`#112\relax}
%    \end{macrocode}
% \end{macro}
%
%
% \begin{macro}{\@vobeyspaces}
% \changes{v1.5}{1993/10/11}{Changed definition to not use \cs{gdef}.}
% \changes{v1.1a}{1989/10/16}{\cs{@vobeyspaces} added.}
%    宏 |\@vobeyspaces| 使输入中的空格在输出中打印为空格。
%    \begin{macrocode}
\begingroup
\catcode`\ =\active%
%    \end{macrocode}
%    因为空格是活动字符，我们无法把下面的代码整齐缩进——
%    否则行首的空格会被当作 \cs{@ifl@t@r} 的第三个和第四个参数。
% \changes{v1.5v}{2023/11/06}{\cs{@vobeytabs} added when available (gh/1160)}
%    \begin{macrocode}
\@ifl@t@r\fmtversion{2023-11-01}%
{\def\x{\def\@vobeyspaces{\catcode`\ \active\let \@xobeysp\@vobeytabs}}}%
{\def\x{\def\@vobeyspaces{\catcode`\ \active\let \@xobeysp}}}%
\expandafter\endgroup\x
%    \end{macrocode}
% \end{macro}
%
%
% \begin{macro}{\@xobeysp}
% \changes{v1.1a}{1989/10/16}{\cs{@xobeysp} added.}
%    宏 |\@xobeysp| 在输出中产生恰好一个空格，
%    并防止在它之前断行。（|\@M| 是数字 $10000$ 的缩写。）
%    \begin{macrocode}
\def\@xobeysp{\leavevmode\penalty\@M\ }
%    \end{macrocode}
% \end{macro}
%
%
% \begin{macro}{\verbatim@line}
% \changes{v1.2d}{1989/11/29}{Introduced token register
%                \cs{verbatim@line}.}
%    我们使用一个新定义的记号寄存器 |\verbatim@line|
%    作为字符缓冲区。
%    \begin{macrocode}
\newtoks\verbatim@line
%    \end{macrocode}
% \end{macro}
%
% 下面四个宏以适合 \texttt{verbatim} 和 \texttt{verbatim*} 环境的方式
% 全局定义。
% \begin{macro}{\verbatim@startline}
% \begin{macro}{\verbatim@addtoline}
% \begin{macro}{\verbatim@processline}
%    |\verbatim@startline| 通过清空字符缓冲区（|\verbatim@line|）
%    来开始处理一行。
%    \begin{macrocode}
\def\verbatim@startline{\verbatim@line{}}
%    \end{macrocode}
%    |\verbatim@addtoline| 把其参数中的记号添加到缓冲区寄存器
%    |\verbatim@line|，不展开它们。
%    \begin{macrocode}
\def\verbatim@addtoline#1{%
  \verbatim@line\expandafter{\the\verbatim@line#1}}
%    \end{macrocode}
%    在 \texttt{verbatim} 或 \texttt{verbatim*} 环境中处理一行
%    意味着打印它。
% \changes{v1.2c}{1989/10/31}{Changed \cs{@@par} to \cs{par} in
%    \cs{verbatim@processline}.  Removed \cs{leavevmode} and \cs{null}
%    (i.e.\ the empty \cs{hbox}).}
%    结束一行意味着我们必须开始新段落。我们用 |\par| 来做这件事。
%    注意 |\par| 在 |\@verbatim| 中被重定义，以迫使 \TeX{} 进入水平模式
%    并插入一个空盒子，这样输入中的空行才会在输出中出现。
% \changes{v1.2f}{1990/01/31}{\cs{verbatim@startline} removed.}
%    \begin{macrocode}
\def\verbatim@processline{\the\verbatim@line\par}
%    \end{macrocode}
% \end{macro}
% \end{macro}
% \end{macro}
%
% \begin{macro}{\verbatim@finish}
%    默认情况下，|\verbatim@finish| 处理剩余字符。
%    调用这个宏时我们面临如下问题：当遇到 |\end{verbatim}| 命令时，
%    |\verbatim@processline| 被调用来处理该命令同一行上前面的字符。
%    如果没有字符，而我们不检查这种情况，就会输出一个空行。
%
%    如果行为空，|\the\verbatim@line| 展开为空。为检测这一点，
%    我们使用一个与 \TeX{}book 第 376 页类似的技巧，
%    但用 |$|\ldots|$| 代替那里的 |!| 记号。这些 |$| 记号的类别码
%    永远不可能与记号寄存器 |\verbatim@line| 中可能出现的 |$| 记号
%    相同，因为后者总是以类别码 $12$（其他）读取的。
%    注意 |\ifcat| 会展开后面的记号，使 |\the\verbatim@line|
%    被累积的字符替换。
% \changes{v1.2d}{1989/11/29}{Changed \cs{ifx} to \cs{ifcat} test.}
% \changes{v1.1b}{1989/10/18}{Corrected bug in if test (found by CRo).}
%    \begin{macrocode}
\def\verbatim@finish{\ifcat$\the\verbatim@line$\else
  \verbatim@processline\fi}
%    \end{macrocode}
% \end{macro}
%
%
% \subsection{\texttt{verbatim} 与 \texttt{verbatim*} 环境}
%
% \begin{macro}{\verbatim@font}
% \changes{v1.2f}{1990/01/31}{\cs{@lquote} macro removed.}
% \changes{v1.1b}{1989/10/18}{\cs{@noligs} removed.  Code inserted
%                           directly into \cs{verbatim@font}.}
% \changes{v1.1a}{1989/10/16}{\cs{verbatim@font} added.}
% \changes{v1.1a}{1989/10/16}{\cs{@noligs} added.}
% \changes{v1.1a}{1989/10/16}{\cs{@lquote} added.}
%    我们首先定义宏 |\verbatim@font|，用于选择字体并设置与字体相关的参数。
%    然后展开 |\@noligs|（定义在 \LaTeXe{} 内核中）。它除了做其他事之外，
%    会遍历 |\verbatim@nolig@list| 以避免某些连字。
%    |\verbatim@nolig@list| 是 \LaTeXe{} 内核中定义的宏，展开为
%    \begin{verbatim}
%    \do\`\do\<\do\>\do\,\do\'\do\-
%\end{verbatim}
%    这个列表中的所有字符在某种字体中都可能成为连字的一部分。
% \changes{v1.2f}{1990/01/31}{\cs{@lquote} macro removed.}
% \changes{v1.4c}{1990/10/18}{Added \cs{leavevmode}.}
% \changes{v1.4k}{1992/07/13}{Replaced Blank after $96$ by \cs{relax}.
%                           (Proposed by Dan Dill.)}
% \changes{v1.5}{1993/10/11}{Definition changed according to new code
%          in latex.tex and to avoid global definition.}
% \changes{v1.5c}{1994/02/07}{Changed to use new font switching
%                           commands.}
% \changes{v1.5m}{2000/01/07}{Disable hyphenation even if the font
%    allows it.}
% \changes{v1.5q}{2003/08/22}{Use \cs{@noligs}, as it is by now properly
%    defined in the \LaTeXe{} kernel.}
%    \begin{macrocode}
\def\verbatim@font{\normalfont\ttfamily
                   \hyphenchar\font\m@ne
                   \@noligs}
%    \end{macrocode}
% \end{macro}
%
%
% \begin{macro}{\@verbatim}
% \changes{v1.1a}{1989/10/16}{\cs{@verbatim} added.}
%    宏 |\@verbatim| 负责做好各项设置。首先，插入
%    |\every@verbatim| 钩子的记号。然后启动一个 \texttt{trivlist}
%    环境并插入其第一个 |\item| 命令。\texttt{verbatim} 或
%    \texttt{verbatim*} 环境的每一行都将被当作独立段落处理。
% \changes{v1.2e}{1990/01/15}{Added \cs{every@verbatim} hook.}
% \changes{v1.5b}{1994/01/24}{Removed optional argument of \cs{item}.}
%    \begin{macrocode}
\def\@verbatim{\the\every@verbatim
  \trivlist \item \relax
%    \end{macrocode}
% \changes{v1.5b}{1994/01/24}{Set \texttt{@inlabel} switch to false.}
% \changes{v1.5f}{1994/10/25}{Removed setting of \texttt{@inlabel}
%                             switch again.}
% \changes{v1.3c}{1990/02/26}{Removed extra vertical space.
%           Suggested by Frank Mittelbach.}
% \changes{v1.5h}{1995/09/21}{Added the space again, since it is
%           necessary for correct vertical spacing if \texttt{verbatim}
%           is nested inside \texttt{quote}.}
%    下面这段额外的垂直间距是为了与 \LaTeX{} 内核兼容：否则，
%    使用本 |verbatim| 宏包会改变嵌套在 |quote| 环境中的
%    |verbatim| 环境的垂直间距。
%    \begin{macrocode}
  \if@minipage\else\vskip\parskip\fi
%    \end{macrocode}
% \changes{v1.4k}{1992/07/13}{Added setting for
%        \cs{@beginparpenalty}. Suggested by Frank Mittelbach.}
%    段落参数被适当设置：环境开头的惩罚值、左右边距、段落缩进、
%    填充最后一行的胶值，以及段落之间的垂直间距。后者必须为零，
%    因为我们不希望在行与行之间添加额外间距。
%    \begin{macrocode}
  \@beginparpenalty \predisplaypenalty
  \leftskip\@totalleftmargin\rightskip\z@
  \parindent\z@\parfillskip\@flushglue\parskip\z@
%    \end{macrocode}
% \changes{v1.1b}{1989/10/18}{Added resetting of \cs{parshape}
%                           if at beginning of a list.
%                           (Problem pointed out by Chris Rowley.)}
%    这里有一点要说明：\texttt{list} 环境使用 \TeX{} 的 |\parshape|
%    原语来为列表第一行设置特殊缩进。如果列表以 \texttt{verbatim}
%    环境开头，这个 |\parshape| 仍然生效。因此我们必须显式重置这个
%    内部参数。我们可以通过把 $0$ 赋给 |\parshape| 来做到。
%    不过有一个更简单的办法：直接告诉 \TeX{} 开始新段落。
%    正如 \TeX{}book 第 103 页所述，这会把 |\parshape| 重置为零。
% \changes{v1.1c}{1989/10/19}{Replaced explicit resetting of
%                           \cs{parshape} by \cs{@@par}.}
%    \begin{macrocode}
  \@@par
%    \end{macrocode}
%    现在我们确保 |\par| 有正确的定义，即迫使 \TeX{} 进入水平模式
%    并包含一个空盒子。这是为了确保空行确实在输出中出现。
%    之后我们插入 |\interlinepenalty|，因为 \TeX{} 不会主动在段落
%    （这里是行）之间添加惩罚值。否则即使存在 |\samepage| 声明，
%    |verbatim| 环境也可能被跨页拆开。
%
%    不过，在顶对齐的 minipage 中，这会导致顶部多出一个空行。
%    因此需要一个稍复杂的结构。这里重要的一点是把 |\leavevmode|
%    作为第一行的第一个宏，例如当列表项的第一个内容就是一个空
%    verbatim 行时。
% \changes{v1.2c}{1989/10/31}{Definition of \cs{par} added.
%                           Ensures identical behaviour for
%                           verbatim and \cs{verbatiminput}.
%                           Problem pointed out by Chris.}
% \changes{v1.4d}{1991/04/24}{\cs{penalty}\cs{interlinepenalty} added.
%                           Necessary to avoid page breaks in
%                           the scope of a \cs{samepage} declaration.}
% \changes{v1.5b}{1994/01/24}{Improved definition of \cs{par} to work
%                           under all circumstances.}
% \changes{v1.5f}{1994/10/25}{\cs{leavevmode} added for first line.}
%    \begin{macrocode}
  \def\par{%
    \if@tempswa
      \leavevmode\null\@@par\penalty\interlinepenalty
    \else
      \@tempswatrue
      \ifhmode\@@par\penalty\interlinepenalty\fi
    \fi}%
%    \end{macrocode}
%    为了在环境不包含任何文本时避免错误信息，我们重定义
%    |\@noitemerr|，它在这种情况下会被 |\endtrivlist| 调用。
% \changes{v1.4j}{1992/06/30}{Introduced warning instead of error
%        for empty body of verbatim text.
%        Suggested by Nelson Beebe.}
%    \begin{macrocode}
  \def\@noitemerr{\@warning{No verbatim text}}%
%    \end{macrocode}
%    现在调用 |\obeylines| 使行尾字符成为活动字符。
%    \begin{macrocode}
  \obeylines
%    \end{macrocode}
%    所有特殊字符的类别码都改为 $12$（其他）。
%    \changes{v1.5i}{1996/06/04}{Moved \cs{verbatim@font} after
%           \cs{dospecials}.}
%    \begin{macrocode}
  \let\do\@makeother \dospecials
%    \end{macrocode}
%    我们切换到要使用的字体。
%    \begin{macrocode}
  \verbatim@font
%    \end{macrocode}
%    为了避免在标签盒子之后断行，我们去掉列表宏在那里放置的惩罚值：
%    又一次用到了 |\unpenalty|！
% \changes{v1.5f}{1994/10/25}{Change to \cs{everypar} added.}
%    \begin{macrocode}
  \everypar \expandafter{\the\everypar \unpenalty}}
%    \end{macrocode}
% \end{macro}
%
%
% \begin{macro}{\verbatim}
% \begin{macro}{\verbatim*}
%    现在定义顶层宏。|\verbatim| 略有改动：做好设置后调用
%    |\verbatim@start|。
% \changes{v1.5l}{1999/12/14}{Added \cs{begingroup} for cases where
%    \cs{verbatim} is used directly, rather than in \cs{begin}: see
%    pr/3115.}
%    这一切都在一个组内完成，这样 |\verbatim| 就可以不通过 |\begin|
%    而直接使用。
%    \begin{macrocode}
\def\verbatim{\begingroup\@verbatim \frenchspacing\@vobeyspaces
              \verbatim@start}
%    \end{macrocode}
%    |\verbatim*| 按类似方式定义。
%    \begin{macrocode}
\@namedef{verbatim*}{\begingroup\@verbatim
%    \end{macrocode}
%
% \changes{v1.5r}{2019/11/10}{Support kernel \cs{verbvisiblespace} (gh/212)}
%    \begin{macrocode}
  \@setupverbvisiblespace\@vobeyspaces
  \verbatim@start}
%    \end{macrocode}
% \end{macro}
% \end{macro}
%
% \begin{macro}{\endverbatim}
% \begin{macro}{\endverbatim*}
%    要结束 \texttt{verbatim} 和 \texttt{verbatim*} 环境，
%    只需结束 |\@verbatim| 中启动的 \texttt{trivlist} 环境并关闭相应的组。
% \changes{v1.5l}{1999/12/14}{Added \cs{endgroup} for cases where
%    \cs{endverbatim} is used directly, rather than in \cs{end}: see
%    pr/3115.}
% \changes{v1.5n}{2000/08/03}{Added \cs{@endpetrue}: needed when
%    faking such a \cs{end} (pr/3234).}
% \changes{v1.5o}{2000/08/23}{Changed \cs{@endpetrue} to \cs{@doendpe}:
%    see (pr/3234).}
%    \begin{macrocode}
\def\endverbatim{\endtrivlist\endgroup\@doendpe}
\expandafter\let\csname endverbatim*\endcsname =\endverbatim
%    \end{macrocode}
% \end{macro}
% \end{macro}
%
%
% \subsection{\texttt{comment} 环境}
%
% \begin{macro}{\comment}
% \begin{macro}{\endcomment}
% \changes{v1.1c}{1989/10/19}{Added \cs{@bsphack}/\cs{@esphack} to the
%            \texttt{comment} environment.  Suggested by Chris Rowley.}
%    |\comment| 宏与 |\verbatim*| 类似。不过我们不需要切换字体或设置
%    诸如 |\parindent|、|\parskip| 之类的特殊排版参数。
%    只需把所有特殊字符的类别码设为 $12$（其他），
%    把 |^^M|（行尾字符）的类别码设为 $13$（活动）。
%    后者是下面内部宏进行宏参数定界匹配所需要的。
%    与 |\verbatim|、|\verbatim*| 宏的默认定义不同，
%    我们把 |\verbatim@addtoline| 定义为丢弃其参数，
%    把 |\verbatim@processline|、|\verbatim@startline|、
%    |\verbatim@finish| 定义为空操作。然后调用 |\verbatim@|。
%    但我们做的第一件事是调用 |\@bsphack|，
%    使该环境对间距完全没有影响。
% \changes{v1.1c}{1989/10/19}{Changed \cs{verbatim@start} to
%                           \cs{verbatim@}.  Suggested by Chris Rowley.}
% \changes{v1.1c}{1989/10/19}{\cs{verbatim@startline} and
%                           \cs{verbatim@finish} are now
%                           also redefined to do nothing.}
%    \begin{macrocode}
\def\comment{\@bsphack
             \let\do\@makeother\dospecials\catcode`\^^M\active
             \let\verbatim@startline\relax
             \let\verbatim@addtoline\@gobble
             \let\verbatim@processline\relax
             \let\verbatim@finish\relax
             \verbatim@}
%    \end{macrocode}
%    |\endcomment| 非常简单：它只调用 |\@esphack| 来处理间距。
%    |\end| 宏关闭组，因此负责恢复我们所做的一切更改。
%    \begin{macrocode}
\let\endcomment=\@esphack
%    \end{macrocode}
% \end{macro}
% \end{macro}
%
%
%
% \subsection{主循环}
%
% 下面是最巧妙的部分：在定义这些宏时，我们不仅需要以正常类别码
% 使用特殊字符 |\|、|{| 和 |}|，还需要以类别码 $12$（其他）使用它们。
% 我们通过下面的技巧实现：先告诉 \TeX{}，|\|、|{| 和 |}|
% 分别是 |!|、|[| 和 |]| 的小写形式。然后把每个应该以类别码 $12$
% 读取的 |\|、|{|、|}| 分别替换为 |!|、|[|、|]|，
% 把整串记号交给 |\lowercase|，因为我们知道这个原语不会改变类别码！
%
% 但首先我们必须确保 |!|、|[|、|]| 本身有正确的类别码！
% \changes{v1.3b}{1990/02/07}{Introduced \cs{vrb@catcodes} instead
%                  of explicit setting of category codes.}
% 为了允许对这些类别码做特殊设置，我们把它们的设置隐藏在宏
% |\vrb@catcodes| 中。如果它已经定义，就跳过我们的新定义。
%    \begin{macrocode}
\@ifundefined{vrb@catcodes}%
  {\def\vrb@catcodes{%
     \catcode`\!12\catcode`\[12\catcode`\]12}}{}
%    \end{macrocode}
% 这个技巧使我们的代码可以用于其他类别码生效的场合。
%
% 我们启动一个组，使类别码改动局部化。
%    \begin{macrocode}
\begingroup
 \vrb@catcodes
 \lccode`\!=`\\ \lccode`\[=`\{ \lccode`\]=`\}
%    \end{macrocode}
% \changes{v1.2f}{1990/01/31}{Code for TABs removed.}
%    我们还需要把行尾字符 |^^M| 作为活动字符。
%    如果直接写 |\catcode`\^^M=\active|，那么在下面每个宏定义行的末尾
%    都会出现一个不想要的活动行尾字符。因此我们用与上面相同的技巧：
%    用波浪号 |~| 代替 |^^M|，假装后者是前者的小写形式。
%    因此我们现在必须确保波浪号字符有类别码 $13$（活动）。
%    \begin{macrocode}
 \catcode`\~=\active \lccode`\~=`\^^M
%    \end{macrocode}
%    使用 |\lowercase| 原语带来一个问题：下面的代码中需要使用大写字母
%    |C|，并且必须保持其大小写。因此我们添加如下命令：
%    \begin{macrocode}
 \lccode`\C=`\C
%    \end{macrocode}
%    现在开始交给 |\lowercase| 的记号表。我们用下面这个小技巧
%    （由 Bernd Raichle 提出）：交给 |\lowercase| 的记号表中第一个记号
%    是 |\endgroup| 原语。这意味着 \TeX{} 在 |\lowercase| 完成操作后
%    立即处理它，从而结束上面 |\begingroup| 启动的组。
%    这样就避免了所有宏的全局定义。
%    \begin{macrocode}
 \lowercase{\endgroup
%    \end{macrocode}
% \begin{macro}{\verbatim@start}
%    |\verbatim@start| 的作用是检查 |\begin{verbatim}| 同一行上
%    是否还有其他字符，并把它们当作单独一行处理。另一方面，如果当前行
%    上没有剩余字符，我们只会遇到一个行尾字符。|\verbatim@start|
%    通过先抓取后面的字符（即其参数）来完成任务。然后把这个参数与
%    活动的行尾字符 |^^M| 比较。
%    \begin{macrocode}
    \def\verbatim@start#1{%
      \verbatim@startline
      \if\noexpand#1\noexpand~%
%    \end{macrocode}
%    如果为真，我们把控制权交给 |\verbatim@| 来处理下一行。
%    我们用 |\next| 作为继续工作的宏。
%    \begin{macrocode}
        \let\next\verbatim@
%    \end{macrocode}
%    否则，我们把 |\next| 定义为展开为对 |\verbatim@| 的调用，
%    后面紧跟刚读入的字符，使它被重新插回文本。这意味着该行上剩余的
%    字符被当作单独一行处理。
%    \begin{macrocode}
      \else \def\next{\verbatim@#1}\fi
%    \end{macrocode}
%    Finally we call |\next|.
%    \begin{macrocode}
      \next}%
%    \end{macrocode}
% \end{macro}
%
% \begin{macro}{\verbatim@}
%    The three macros |\verbatim@|, |\verbatim@@|,
%    三个宏 |\verbatim@|、|\verbatim@@|、|\verbatim@@@| 构成
%    \texttt{verbatim} 环境的"主循环"。|\verbatim@| 的作用是恰好读取
%    一行输入。|\verbatim@@| 和 |\verbatim@@@| 协同工作，
%    找出该行中是否出现四个字符 |\end|（均为类别码 $12$（其他））。
%    如果出现，|\verbatim@@@| 会调用 |\verbatim@test| 检查这个 |\end|
%    是否是 |\end{verbatim}| 的一部分，如果是就结束环境。
%    否则我们就像什么都没发生一样继续。下面看 |\verbatim@| 的定义：
% \changes{v1.1a}{1989/10/16}{Replaced \cs{verbatim@@@} by \cs{@nil}.}
%    \begin{macrocode}
    \def\verbatim@#1~{\verbatim@@#1!end\@nil}%
%    \end{macrocode}
%    注意：在这个宏的定义真正生效之前，|!| 字符已经被管辖本代码的
%    |\lowercase| 原语替换为一个类别码 $12$（其他）的 |\|。
%    也就是说，它取一行，在后面放上 |\end|（四个字符记号）和
%    |\@nil|（一个控制序列记号）作为定界符，然后调用 |\verbatim@@|。
% \end{macro}
%
% \begin{macro}{\verbatim@@}
%    |\verbatim@@| 把直到下一次出现四个字符 |\end| 之前的所有内容
%    作为其参数。
%    \begin{macrocode}
    \def\verbatim@@#1!end{%
%    \end{macrocode}
%    也就是说：如果原行中没有出现它们，那么参数 |#1| 就是整行输入，
%    下一个要处理的记号是 |\@nil|。但是，如果四个字符 |\end|
%    是原行的一部分，那么 |#1| 由 |\end| 前面的字符组成，
%    下一个记号是后面的字符（永远记住该行被加长了五个记号）。
%    无论 |#1| 是什么，它都是 verbatim 文本，所以 |#1| 被添加到
%    当前正在构建的行中。
%    \begin{macrocode}
       \verbatim@addtoline{#1}%
%    \end{macrocode}
%    输入流中的下一个记号对我们特别重要。因此在调用 |\verbatim@@@|
%    之前，|\futurelet| 把 |\next| 定义为等于该记号。
%    \begin{macrocode}
       \futurelet\next\verbatim@@@}%
%    \end{macrocode}
% \end{macro}
%
% \begin{macro}{\verbatim@@@}
% \changes{v1.1a}{1989/10/16}{Replaced \cs{verbatim@@@} by
%                           \cs{@nil} where used as delimiter.}
%    |\verbatim@@@| 现在读取当前行上剩余的记号，直到最后的 |\@nil| 记号。
%    \begin{macrocode}
    \def\verbatim@@@#1\@nil{%
%    \end{macrocode}
%    如果上述两种情况中的第一种发生，即该行上没有 |\end| 字符，
%    那么 |#1| 为空，且 |\next| 等于 |\@nil|。这很容易检查。
%    \begin{macrocode}
       \ifx\next\@nil
%    \end{macrocode}
%    如果为真，这就是一个简单行。我们处理迄今累积的行来结束它。
%    然后准备读取下一行。
% \changes{v1.2f}{1990/01/31}{Added \cs{verbatim@startline}.}
%    \begin{macrocode}
         \verbatim@processline
         \verbatim@startline
         \let\next\verbatim@
%    \end{macrocode}
%    否则我们必须检查这些 |\end| 记号后面跟着什么。
%    \begin{macrocode}
       \else
%    \end{macrocode}
%    在继续之前，最好停下来想一下我们在哪儿：我们刚读完四个字符记号
%    |\end|，现在必须检查后面是否跟着环境名（用花括号括起来）。
%    为此我们定义一个叫 |\@tempa| 的宏，它恰好读取一个字符并决定
%    下一步做什么。这个宏应该做下面的事：跳过空格，直到遇到左花括号
%    或行尾。但重要的是记住跳过了哪些字符。|\end|\meta{可选空格}|{|
%    这些字符可能是 verbatim 文本的一部分，即这些字符必须被打印。
%
%    例如，假设当前行包含
%    \begin{verbatim*}
%      \end {AVeryLongEnvironmentName}
%\end{verbatim*}
%    正如我们很快会看到的，这里实现的扫描机制在读完右花括号之前
%    不会发现这是要打印的文本。因此我们需要一种累积已读字符的方法，
%    以便在必要时把它们重新插回。记号寄存器 |\@temptokena| 用于此目的。
%
%    在做这件事之前，我们必须先去掉行末多余的 |\end| 记号。
% \changes{v1.4j}{1992/06/30}{Removed use of \cs{toks@}. Suggested by
%           Bernd Raichle.}
%    为此我们定义一个临时宏，其参数由 |\end\@nil|（四个字符记号和一个
%    控制序列记号）界定，用于下面处理该行剩余部分——先在该行末尾追加一个
%    |\@nil| 记号。（注意这个记号永远不可能出现在 |#1| 中。）
%    我们用下面的 |\@tempa| 定义来获取该行（第一个 |\end| 之后）的剩余部分。
%    \begin{macrocode}
         \def\@tempa##1!end\@nil{##1}%
%    \end{macrocode}
%    我们前面已经提到，用记号寄存器 |\@temptokena| 来记住跳过的字符，
%    以防再次需要它们。我们用在 |\@tempa| 调用中丢弃的那个 |\end|
%    来初始化它。
%    \begin{macrocode}
         \@temptokena{!end}%
%    \end{macrocode}
%    现在我们要调用 |\verbatim@test| 来处理当前行上剩余的字符。
%    但等一下：我们不能直接调用这个宏，因为我们已经读完了整行。
%    因此我们必须先展开宏 |\@tempa|，把这些字符重新插回到
%    |\verbatim@test| 记号之后。末尾追加一个 |^^M| 字符表示行尾。
%    （记住这个字符伪装成波浪号。）
% \changes{v1.2}{1989/10/20}{Taken local definition of \cs{@tempa} out
%                          of \cs{verbatim@@@} and introduced
%                          \cs{verbatim@test} instead.}
%    \begin{macrocode}
         \def\next{\expandafter\verbatim@test\@tempa#1\@nil~}%
%    \end{macrocode}
%    差不多就这些了，但我们还需要调用 |\next| 来完成工作。
%    \begin{macrocode}
       \fi \next}%
%    \end{macrocode}
% \end{macro}
%
%
% \begin{macro}{\verbatim@test}
% \changes{v1.2}{1989/10/20}{Introduced \cs{verbatim@test}.}
%    我们定义 |\verbatim@test| 来依次检查每个记号。
%    \begin{macrocode}
    \def\verbatim@test#1{%
%    \end{macrocode}
%    首先，我们把 |\next| 设为 |\verbatim@test|，
%    以防这个宏需要递归调用来跳过空格。
%    \begin{macrocode}
           \let\next\verbatim@test
%    \end{macrocode}
%    我们必须区分四种情况：
%    \begin{enumerate}
%      \item 下一个记号是 |^^M|，即我们到达了行尾。这意味着没有发现
%            任何特殊内容。注意我们在下面的比较中使用 |\if|，
%            这样字符的类别码就无关紧要了。
%    \begin{macrocode}
           \if\noexpand#1\noexpand~%
%    \end{macrocode}
%            我们把记号寄存器 |\@temptokena| 中累积的字符添加到当前行。
%            由于 |\verbatim@addtoline| 不展开其参数，我们必须在此处展开。
%            然后用 |\let| 把 |\next| 设为 |\verbatim@|，
%            准备读取下一行。
% \changes{v1.2f}{1990/01/31}{Added \cs{verbatim@startline}.}
%    \begin{macrocode}
             \expandafter\verbatim@addtoline
               \expandafter{\the\@temptokena}%
             \verbatim@processline
             \verbatim@startline
             \let\next\verbatim@
%    \end{macrocode}
%      \item 后面是空格字符。这是允许的，所以我们把它添加到
%            |\@temptokena| 并继续。
%    \begin{macrocode}
           \else \if\noexpand#1
             \@temptokena\expandafter{\the\@temptokena#1}%
%    \end{macrocode}
% \changes{v1.2f}{1990/01/31}{Code for TABs removed.}
%      \item 后面是左花括号。这是最有趣的情况。我们现在必须收集字符，
%            直到读到右花括号，并检查它们是否构成环境名。
%            这将由 |\verbatim@testend| 完成，所以这里我们把
%            |\next| 设为这个宏。我们将再次逐字符处理该行剩余部分。
%            构成环境名的字符将累积在 |\@tempc| 中。
%            我们把这个宏初始化为展开为空。
% \changes{v1.3b}{1990/02/07}{\cs{noexpand} added.}
%    \begin{macrocode}
           \else \if\noexpand#1\noexpand[%
             \let\@tempc\@empty
             \let\next\verbatim@testend
%    \end{macrocode}
%            注意：当这个宏被定义时，|[| 字符将是一个 |{|。
%      \item 任何其他字符都意味着这个 |\end| 是 verbatim 文本的一部分。
%            把这些字符添加到当前行，并准备调用 |\verbatim@|
%            来处理该行剩余部分。
%  \changes{v1.0f}{1989/10/09}{Fixed \cs{end} \cs{end} bug
%                            found by Chris Rowley}
%    \begin{macrocode}
           \else
             \expandafter\verbatim@addtoline
               \expandafter{\the\@temptokena}%
             \def\next{\verbatim@#1}%
           \fi\fi\fi
%    \end{macrocode}
%    \end{enumerate}
%    这个宏做的最后一件事是调用 |\next| 来继续处理。
%    \begin{macrocode}
           \next}%
%    \end{macrocode}
% \end{macro}
%
% \begin{macro}{\verbatim@testend}
%    当看到 |\end|\meta{可选空格}|{| 时调用 |\verbatim@testend|。
%    它的任务是扫描直到下一个 |}| 的所有内容，并调用 |\verbatim@@testend|。
%    如果没找到 |}|，它必须重新插回所读字符并返回 |\verbatim@|。
%    下面的定义与 |\verbatim@test| 类似：它取下一个字符并决定做什么。
% \changes{v1.2}{1989/10/20}{Removed local definition of \cs{@tempa}
%                          from \cs{verbatim@testend} which now
%                          does the work itself.}
%    \begin{macrocode}
    \def\verbatim@testend#1{%
%    \end{macrocode}
%    同样，我们有四种情况：
%    \begin{enumerate}
%      \item |^^M|：由于当前行中没有找到 |}|，把这些字符添加到缓冲区。
%            为避免展开 |\@temptokena| 和 |\@tempc| 的复杂构造，
%            我们分两步做。然后继续用 |\verbatim@| 处理下一行。
% \changes{v1.2f}{1990/01/31}{Added \cs{verbatim@startline}.}
%    \begin{macrocode}
         \if\noexpand#1\noexpand~%
           \expandafter\verbatim@addtoline
             \expandafter{\the\@temptokena[}%
           \expandafter\verbatim@addtoline
             \expandafter{\@tempc}%
           \verbatim@processline
           \verbatim@startline
           \let\next\verbatim@
%    \end{macrocode}
%      \item |}|：调用 |\verbatim@@testend| 检查这是否是正确的环境名。
% \changes{v1.3b}{1990/02/07}{\cs{noexpand} added.}
%    \begin{macrocode}
         \else\if\noexpand#1\noexpand]%
           \let\next\verbatim@@testend
%    \end{macrocode}
%  \changes{v1.0f}{1989/10/09}{Introduced check for
%              {\tt\string\verb!\string|!\string\!\string|} to fix
%              single brace bug found by Chris Rowley}
%      \item |\|：这个字符不应出现在环境名中。因此我们停止收集字符。
%            原则上，同样的理由也适用于其他字符，例如 |{|。
%            不过 |\| 是一种特殊情况，因为它可能是 |\end| 的第一个字符。
%            这意味着我们必须重新寻找 |\end{|\meta{环境名}|}|。
%            注意我们在 |!| 前面加了一个 |\noexpand| 原语，
%            以防止它是一个活动字符。
% \changes{v1.3b}{1990/02/07}{\cs{noexpand} added.}
%    \begin{macrocode}
         \else\if\noexpand#1\noexpand!%
           \expandafter\verbatim@addtoline
             \expandafter{\the\@temptokena[}%
           \expandafter\verbatim@addtoline
             \expandafter{\@tempc}%
           \def\next{\verbatim@!}%
%    \end{macrocode}
%      \item 任何其他字符：收集它并继续。我们不能用 |\edef| 来定义
%            |\@tempc|，因为它的替换文本可能包含活动字符记号。
%    \begin{macrocode}
         \else \expandafter\def\expandafter\@tempc\expandafter
           {\@tempc#1}\fi\fi\fi
%    \end{macrocode}
%    \end{enumerate}
%    和前面一样，这个宏最后通过递归调用来处理下一个字符（如果合适的话）。
%    \begin{macrocode}
         \next}%
%    \end{macrocode}
% \end{macro}
%
% \begin{macro}{\verbatim@@testend}
%    与前面的宏不同，|\verbatim@@testend| 很简单：
%    它只需检查 |\begin{|\ldots|}| 是否与对应的
%    |\end{|\ldots|}| 匹配。
%    \begin{macrocode}
    \def\verbatim@@testend{%
%    \end{macrocode}
%    我们再次用 |\next| 来定义要做的事。记住当前环境名存放在
%    |\@currenvir| 中，由 |\verbatim@testend| 累积的字符在 |\@tempc| 中。
%    所以我们只需比较这两者，如果匹配就准备执行
%    |\end{|\meta{当前环境}|}| 宏。在做这件事之前，我们调用
%    |\verbatim@finish| 来处理最后一行。我们通过 |\edef| 定义 |\next|，
%    这样 |\@currenvir| 就会被其展开替换。因此我们需要 |\noexpand|
%    来抑制此时 |\end| 的展开。
%    \begin{macrocode}
       \ifx\@tempc\@currenvir
         \verbatim@finish
         \edef\next{\noexpand\end{\@currenvir}%
%    \end{macrocode}
%    没有这个技巧，|\end| 命令将无法正确检查其参数是否与当前环境名匹配，
%    你会得到一条有趣的 \LaTeX{} 错误信息，例如：
%    \begin{verbatim}
%! \begin{verbatim*} ended by \end{verbatim*}.
%\end{verbatim}
%    但是，那一行上剩余的字符怎么办？我们调用 |\verbatim@rescan|
%    来处理。它的第一个参数是刚结束的环境名（以防我们还需要它）。
%    |\verbatim@rescan| 把待重新处理的字符列表作为第二个参数。
%    （这个记号表是由 |\verbatim@@@| 插入到当前宏之后的。）
%    由于我们仍在 |\edef| 中，我们用 |\noexpand| 保护它。
%    \begin{macrocode}
                    \noexpand\verbatim@rescan{\@currenvir}}%
%    \end{macrocode}
%    如果名称不匹配，我们把迄今读到的所有字符重新插回，并准备调用
%    |\verbatim@| 来处理该行剩余部分。
%    \begin{macrocode}
       \else
         \expandafter\verbatim@addtoline
           \expandafter{\the\@temptokena[}%
           \expandafter\verbatim@addtoline
             \expandafter{\@tempc]}%
         \let\next\verbatim@
       \fi
%    \end{macrocode}
%    最后我们调用 |\next|。
%    \begin{macrocode}
       \next}%
%    \end{macrocode}
% \end{macro}
%
% \begin{macro}{\verbatim@rescan}
%    原则上，|\verbatim@rescan| 可以用来分析 |\end{...}| 命令之后
%    剩余的字符，并假装这些字符是在"标准"类别码生效下被"正确"读取的。
% \footnote{记住它们都是以类别码 $11$（字母）和 $12$（其他）读取的，
%          这样控制序列就不会被识别为控制序列。}
%    但这并非总是可行（当该行剩余部分中有不匹配的花括号时）。
%    此外，我们认为这不值得费力：在 \texttt{verbatim} 或
%    \texttt{verbatim*} 环境之后，输出反正会开始新行，
%    而 |\end{comment}| 很容易单独放在一行上。因此这里没有理由
%    出现任何文本。对于确实在那里放了东西（也许是注释）的用户，
%    我们简单地发出一条警告并丢弃它们。
%    测试方法在 \TeX{}book 附录 D 第 376 页有说明。
%    我们用 |^^M| 代替那里用的 |!| 字符，因为这个字符不可能出现在
%    |#1| 中。两个 |\noexpand| 原语是必要的，以避免活动字符和宏的展开。
%
%    这里还应注意一个额外的微妙之处：记住我们当前正在构建的记号表，
%    会先被 |\lowercase| 原语处理，然后 \TeX{} 才执行这些定义。
%    这意味着 |\@warning| 宏参数中的 |C| 字符必须受到保护，
%    防止被改成 |c|。这就是我们在上面添加 |\lccode`\C=`\C| 赋值的原因。
%    现在我们可以结束 |\lowercase| 的参数以及改变类别码的那个组了。
%    \begin{macrocode}
    \def\verbatim@rescan#1#2~{\if\noexpand~\noexpand#2~\else
        \@warning{Characters dropped after `\string\end{#1}'}\fi}}
%    \end{macrocode}
% \end{macro}
%
% \subsection{\cs{verbatiminput} 命令}
%
% \begin{macro}{\verbatim@in@stream}
%    我们首先分配一个输入流（从可用的 16 个输入流中）。
%\iffalse
%  Vorstellbar ist auch der Aufruf von |`verbatiminput| innerhalb eines
%  |`verbatiminput| (z.B: wenn man |`input|-Anweisungen im zu lesenden
%  File hat und auch diese Files automatisch lesen will).  Dies kann
%  man jedoch nur ermoeglichen, wenn man einen besseren Mechanismus
%  verwendet als es das simple, statische |`newread| darstellt.
%  Vorstellbar fuer eine neuere \LaTeX-Version ist eine (lokale)
%  Allokation des Streams durch ein |`open| und eine Freigabe des
%  Streams durch |`close| oder Verlassen der Gruppe.
%\fi
%    \begin{macrocode}
\newread\verbatim@in@stream
%    \end{macrocode}
% \end{macro}
%
% \begin{macro}{\verbatim@readfile}
%    宏 |\verbatim@readfile| 通过分别调用 |\verbatim@startline| 和
%    |\verbatim@finish| 宏来包围主循环。这样可以确保当文件为空或不存在时，
%    用户仍能正确初始化和结束该命令。\texttt{verbatim} 环境在以空文本
%    调用时有类似行为。
%    \begin{macrocode}
\def\verbatim@readfile#1{%
  \verbatim@startline
%    \end{macrocode}
%    当找不到文件时，我们发出一条警告。
%    \begin{macrocode}
  \openin\verbatim@in@stream #1\relax
  \ifeof\verbatim@in@stream
    \typeout{No file #1.}%
  \else
%    \end{macrocode}
%    此时我们把文件名传给 |\@addtofilelist|，使它出现在
%    |\listfiles| 命令的输出中。
%    此外，我们使用 |\ProvidesFile| 在日志文件中记录一条日志项，
%    并区分通过 |\verbatiminput| 读入的文件与其他文件。
%    \begin{macrocode}
    \@addtofilelist{#1}%
    \ProvidesFile{#1}[(verbatim)]%
%    \end{macrocode}
%    从文件读取时，关闭对行尾字符的识别是有用的。这样我们就不必
%    从行内容中去掉空格。
%    \begin{macrocode}
    \expandafter\endlinechar\expandafter\m@ne
    \expandafter\verbatim@read@file
    \expandafter\endlinechar\the\endlinechar\relax
    \closein\verbatim@in@stream
  \fi
  \verbatim@finish
}
%    \end{macrocode}
% \end{macro}
%
% \begin{macro}{\verbatim@read@file}
%    所有工作都在 |\verbatim@read@file| 中完成。它逐行读取输入文件，
%    并递归调用自己直到文件末尾。
%    \begin{macrocode}
\def\verbatim@read@file{%
  \read\verbatim@in@stream to\next
  \ifeof\verbatim@in@stream
  \else
%    \end{macrocode}
%    对每一行，我们用该行内容调用 |\verbatim@addtoline|。
% \hskip0pt plus 3cm\penalty0\hskip0pt plus -3cm
%    接着调用 |\verbatim@processline|。
%    \begin{macrocode}
    \expandafter\verbatim@addtoline\expandafter{\next}%
    \verbatim@processline
%    \end{macrocode}
%    处理完一行后，我们在读取下一行之前调用 |\verbatim@startline|
%    来重新初始化一切。
%    \begin{macrocode}
    \verbatim@startline
%    \end{macrocode}
%    如果没有 |\expandafter|，每次调用 |\verbatim@read@file| 都会在
%    \TeX 的输入栈中占用空间。\footnote{标准 \TeX{} 在你试图读取
%    超过约 200 行的文件时会报告溢出错误。如果 \textsl{``TeX: The Program''}\S 390
%    中的第一行代码缺失，也会发生同样的错误。}
%    \begin{macrocode}
    \expandafter\verbatim@read@file
  \fi
}
%    \end{macrocode}
% \end{macro}
%
%
% \begin{macro}{\verbatiminput}
%    |\verbatiminput| 首先启动一个组，使字体和类别码改动局部化。
%    然后它根据后面是否跟星号，带附加参数调用宏 |\verbatim@input|。
%    \begin{macrocode}
\def\verbatiminput{\begingroup
%    \end{macrocode}
% 
%   \changes{v1.5x}{2024/01/22}{Added TAB marking support into the
%   starred version (gh/1245)}
%   如果是星号版本，我们标记空格和 TAB，这两个新增部分与 verbatim* 相同。
%    \begin{macrocode}
  \@ifstar{\verbatim@input{\@setupverbvisiblespace\@vobeyspaces}}%
          {\verbatim@input{\frenchspacing\@vobeyspaces}}}
%    \end{macrocode}
% \changes{v1.5z}{2026-08-17}{Add tagging support if format is new.}
%    \begin{macrocode}
\IfDocumentMetadataT
 {
  \def\verbatiminput{\begingroup
    \@ifstar{\verbatim@input{final-code=\legacyverbatimsetup{visible}}}%
            {\verbatim@input{final-code=\legacyverbatimsetup{invisible}}}}
 }
%    \end{macrocode}
% \end{macro}
%
% \begin{macro}{\verbatim@input}
% \changes{1.5k}{1997/04/30}{Have \cs{verbatim@input} check for
%    existence of file.}
%    |\verbatim@input| 首先用标准宏 |\IfFileExists| 检查文件是否存在，
%    它会把找到的文件名留在 |\@filef@und| 中。然后像 |\verbatim| 宏一样
%    做好一切设置。但是，由于 |\@verbatim| 包含对 |\every@verbatim| 的调用，
%    而后者\emph{可能}包含一条 |\input| 语句，会覆盖 |\@filef@und| 的内容，
%    我们需要先展开它来保存。使用 |\@swaptwoargs| 使得
%    |\@filef@und| 的\emph{展开结果}成为 |\verbatim@readfile| 的第二个参数。
% \changes{v1.5t}{2020-07-06}{Expand \cs{@filef@und} before the call
%    of \cs{@verbatim} (gh/222)}
%    \begin{macrocode}
\def\verbatim@input#1#2{%
  \IfFileExists {#2}{%
    \expandafter\@swaptwoargs\expandafter
      {\expandafter{\@filef@und}}%
      {\@verbatim #1\relax
%    \end{macrocode}
%    然后它读入文件，结束由 |\@verbatim| 启动的 \texttt{trivlist} 环境，
%    并关闭组。这会把一切恢复到正常设置。
%    \begin{macrocode}
        \verbatim@readfile}%
      \endtrivlist\endgroup\@doendpe}%
%    \end{macrocode}
%    如果找不到文件，会打印一条多少有点用的信息。最后的 |\endgroup|
%    是为了关闭上面 |\verbatiminput| 中启动的组。
%    \begin{macrocode}
   {\typeout {No file #2.}\endgroup}}
%    \end{macrocode}
% \changes{v1.5z}{2026-08-17}{Add tagging support}
%    \begin{macrocode}
\IfDocumentMetadataT
  {
    \def\verbatim@input#1#2{%
      \IfFileExists {#2}{%
        \expandafter\@swaptwoargs\expandafter
         {\expandafter{\@filef@und}}%
         {%
          \SimpleBlockEnv{verbatim} {#1}\obeyedline       
          \verbatim@readfile}        
          \BlockEnvEnd\endgroup
        }%
     {\typeout {No file #2.}\endgroup}}
  }    
%    \end{macrocode}
%    \begin{macrocode}
%</package>
%    \end{macrocode}
% \end{macro}
%
%
% \subsection{把 verbbatim 文本送入参数}
%
% 实现这一点的一种方法是定义一个宏（命令），其展开结果就是所需的
% verbatim 文本。这个命令随后可以在任何需要 verbatim 文本的地方使用。
% 它可以用在参数中，甚至是移动参数中，但它是脆弱的（至少这里的版本是）。
%
% 下面这段代码声称提供了这个功能。它是我（Chris）约两年前所写东西的
% 大幅修订版。也许还需要进一步修订。
%
% 它只打算作为 |\verb| 的扩展，而不是 \texttt{verbatim} 环境的扩展。
% 因此，它或许应该以对 |\verb| 最合适的方式来处理行尾。
%
% \begin{macro}{\newverbtext}
% 这是产生一个新宏的命令，其展开结果是 verbatim 文本。
% 这个命令本身当然不能用在参数中！它的用法如下：
%
% \begin{verbatim}
%    \newverbtext{\myverb}"^%{ &~_\}}@ #"
% \end{verbatim}
%
% 界定 verbatim 文本的规则与 |\verb| 相同。
%
%    \begin{macrocode}
%<*verbtext>
\def \newverbtext {%
  \@ifstar {\@tempswatrue \@verbtext }{\@tempswafalse \@verbtext *}%
}
%    \end{macrocode}
% \end{macro}
%    我认为在这里使用一个临时开关是安全的：如果不安全，可以用合适的
%    |\let|。
%    \changes{v1.5i}{1996/06/04}{Moved processing of
%              \cs{verbatim@nolig@list} after \cs{dospecials}.}
%    \changes{v1.5p}{2001/03/12}{Added missing right brace in
%               definition of \cs{@verbtext} (PR 3314).}
%    \begin{macrocode}
\def \@verbtext *#1#2{%
   \begingroup
     \let\do\@makeother \dospecials
     \let\do\do@noligs \verbatim@nolig@list
     \@vobeyspaces
     \catcode`#2\active
     \catcode`~\active
     \lccode`\~`#2%
     \lowercase
%    \end{macrocode}
%    我们在这里用一个临时宏和一个技巧，使命令本身的定义可以在组内完成，
%    并且是局部定义（也许有更好的方法来实现这一点）。
%    \begin{macrocode}
     {\def \@tempa ##1~%
           {\whitespaces
%    \end{macrocode}
%    如果这些 |\noexpand| 改成 |\noexpand\protect\noexpand|，
%    会不会让这些东西变得健壮？
%    \begin{macrocode}
            \edef #1{\noexpand \@verbtextmcheck
                     \bgroup
                     \if@tempswa
                       \noexpand \visiblespaces
                     \fi
                     \noexpand \@verbtextsetup
                     ##1%
                     \egroup}%
            }%
      \expandafter\endgroup\@tempa
     }
}
%    \end{macrocode}
%    这为当前模式设置正确类型的组：它绝不能在定义时被展开！
%    \begin{macrocode}
\def \@verbtextmcheck {%
   \relax\ifmmode
           \hbox
         \else
           \leavevmode
           \null
         \fi
}
%    \end{macrocode}
%    这里包含其他一些在定义期间不应展开的东西。
%    \begin{macrocode}
\def \@verbtextsetup {%
   \frenchspacing
   \verbatim@font
   \verbtextstyle
}
%    \end{macrocode}
%    命令 |\verbtextstyle| 是一个文档级钩子，可用于覆盖用
%    |\newverbtext| 命令定义的命令的预定义排版处理。
%
%    |\visiblespaces| 和 |\whitespaces| 是这个钩子可能取值的例子。
%    \begin{macrocode}
\let \verbtextstyle \relax
\def \visiblespaces {\chardef \  32\relax}
\def \whitespaces {\let \ \@@space}
\let \@@space \ %
%</verbtext>
%    \end{macrocode}
%
%
% \section{测试实现}
%
% 为了测试实现并做演示，我们提供了一个额外文件。它可以通过条件
% `\textsf{testdriver}' 抽取。它使用一个单独分发的小输入文件
% `\texttt{verbtest.tst}'。我们再次使用单独的 `+' 守卫。
%    \begin{macrocode}
%<*testdriver>
\documentclass{article}

\usepackage{verbatim}

\newenvironment{myverbatim}%
   {\endgraf\noindent MYVERBATIM:\endgraf\verbatim}%
   {\endverbatim}

\makeatletter

\newenvironment{verbatimlisting}[1]%
 {\def\verbatim@startline{\input{#1}%
    \def\verbatim@startline{\verbatim@line{}}%
    \verbatim@startline}%
  \verbatim}{\endverbatim}

\newwrite\verbatim@out

\newenvironment{verbatimwrite}[1]%
 {\@bsphack
  \immediate\openout \verbatim@out #1
  \let\do\@makeother\dospecials\catcode`\^^M\active
  \def\verbatim@processline{%
    \immediate\write\verbatim@out{\the\verbatim@line}}%
  \verbatim@start}%
 {\immediate\closeout\verbatim@out\@esphack}

\makeatother

\addtolength{\textwidth}{30pt}

\begin{document}

\typeout{}
\typeout{===> Expect ``characters dropped''
         warning messages in this test! <====}
\typeout{}

Text Text Text Text Text Text Text Text Text Text Text
Text Text Text Text Text Text Text Text Text Text Text
Text Text Text Text Text Text Text Text Text Text Text
  \begin{verbatim}
    test
    \end{verbatim*}
    test
    \end{verbatim
    test of ligatures: <`!`?`>
    \endverbatim
    test
    \end  verbatim
    test
    test of end of line:
    \end
    {verbatim}
  \end{verbatim} Further text to be typeset: $\alpha$.
Text Text Text Text Text Text Text Text Text Text Text
Text Text Text Text Text Text Text Text Text Text Text
Text Text Text Text Text Text Text Text Text Text Text
  \begin{verbatim*}
    test
    test
  \end {verbatim*}
Text Text Text Text Text Text Text Text Text Text Text
Text Text Text Text Text Text Text Text Text Text Text
Text Text Text Text Text Text Text Text Text Text Text
  \begin{verbatim*}  bla bla
    test
    test
  \end {verbatim*}
Text Text Text Text Text Text Text Text Text Text Text
Text Text Text Text Text Text Text Text Text Text Text
Text Text Text Text Text Text Text Text Text Text Text
Text Text Text Text Text Text Text Text Text Text Text

First of Chris Rowley's fiendish tests:
\begin{verbatim}
the double end test<text>
\end\end{verbatim}  or even \end \end{verbatim}
%
%not \end\ended??
%\end{verbatim}

Another of Chris' devils:
\begin{verbatim}
the single brace test<text>
\end{not the end\end{verbatim}
%
%not \end}ed??
%\end{verbatim}
Back to my own tests:
  \begin{myverbatim}
    test
    test
  \end {myverbatim} rest of line
Text Text Text Text Text Text Text Text Text Text Text
Text Text Text Text Text Text Text Text Text Text Text
Text Text Text Text Text Text Text Text Text Text Text

Test of empty verbatim:
\begin{verbatim}
\end{verbatim}
Text Text Text Text Text Text Text Text Text Text Text
Text Text Text Text Text Text Text Text Text Text Text
Text Text Text Text Text Text Text Text Text Text Text
  \begin {verbatimlisting}{verbtest.tex}
    Additional verbatim text
      ...
  \end{verbatimlisting}
And here for listing a file:
  \verbatiminput{verbtest.tex}
And again, with explicit spaces:
  \verbatiminput*{verbtest.tex}
Text Text Text Text Text Text Text Text Text Text Text
Text Text Text Text Text Text Text Text Text Text Text
Text Text Text Text Text Text Text Text Text Text Text
  \begin{comment}
    test
    \end{verbatim*}
    test
    \end {comment
    test
    \endverbatim
    test
    \end  verbatim
    test
  \end {comment} Further text to be typeset: $\alpha$.
Text Text Text Text Text Text Text Text Text Text Text
Text Text Text Text Text Text Text Text Text Text Text
Text Text Text Text Text Text Text Text Text Text Text
  \begin{comment}  bla bla
    test
    test
  \end {comment}
Text Text Text Text Text Text Text Text Text Text Text
Text Text Text Text Text Text Text Text Text Text Text
Text Text Text Text Text Text Text Text Text Text Text
Text Text Text Text Text Text Text Text Text Text Text

\begin{verbatimwrite}{verbtest.txt}
asfa<fa<df
sdfsdfasd
asdfa<fsa
\end{verbatimwrite}

\end{document}
%</testdriver>
%    \end{macrocode}
%
%
% \Finale
%
\endinput
%%
